\section{Optimization Problem Formulation}

We formulate a Mixed-Integer Non-Linear Programming (MINLP) problem to jointly optimize the routing decisions, transmit powers, and computational resource allocations. The objective is to minimize a weighted sum of the robots' transmission energy, the system latency, and the drone's energy load. Let $\omega_1, \omega_2, \omega_3$ be user-defined weighting coefficients such that $\omega_1 + \omega_2 + \omega_3 = 1$.


The problem is mathematically formulated as:
\begin{equation}
\min_{\mathbf{x}, \mathbf{P}, \mathbf{f}_D, \mathbf{f}_C} \sum_{i=1}^N \left[ \omega_1 E^{\text{tx}}_i + \omega_2 T_i + \omega_3 E_{\text{Drone}, i} \right]
\end{equation}

\textbf{Subject to:}
\begin{align}
& T_i \le \tau_i, \quad \forall i \in \mathcal{N} \label{eq:c1} \\
& 0 \le P_i \le P^{\text{max}}_{\text{robot}}, \quad \forall i \in \mathcal{N} \label{eq:c2} \\
& \sum_{i=1}^N x_i f_{D,i} \le F^{\text{max}}_{\text{Drone}} \label{eq:c3} \\
& \sum_{i=1}^N (1 - x_i) f_{C,i} \le F^{\text{max}}_{\text{Cloud}} \label{eq:c4} \\
& f_{D,i} \ge 0, \quad f_{C,i} \ge 0, \quad \forall i \in \mathcal{N} \label{eq:c5} \\
& x_i \in \{0, 1\}, \quad \forall i \in \mathcal{N} \label{eq:c6}
\end{align}

Constraint \eqref{eq:c1} ensures that tasks are processed within their maximum tolerable latency limit. Constraint \eqref{eq:c2} bounds the ground robots' transmit power. Constraints \eqref{eq:c3} and \eqref{eq:c4} ensure that the total allocated CPU frequencies do not exceed the respective maximum capacities of the Drone-AP and the Cloud server. Finally, constraints \eqref{eq:c5} and \eqref{eq:c6} define the bounds and domains of the continuous and binary variables, respectively.


% BEGIN PANAMA ADDITIONS: PANAMA under communication delay
\begin{panamaaddition}
\subsection{PANAMA Policy with Cached Mission Information}
We use PANAMA's shared dual-stream action-value network, Double DQN learner,
prioritized experience replay (PER), task reward, and exploration mechanism~\cite{panama}.
Nominal learning parameters follow Table~III of PANAMA, including discount
$\gamma=0.99$, AdamW learning rate $10^{-4}$, minibatch size $M=64$, replay
capacity $2\times10^6$, and target mixing coefficient $\eta_{\rm tar}=0.005$.
The communication mechanism introduces no additional reward, imitation policy,
or recovery action. We consider the task-policy channels for obstacles,
visible robots, and communicated intentions, with UAV-mediated transport.
Terrestrial radio maps and additional SNR-based shaping are outside this
algorithmic extension.

An accepted guidance payload contains an absolute A* route
$\mathcal{P}_i=(\mathbf{w}_{i,0},\ldots,\mathbf{w}_{i,L})$ and reported
distance $\widehat{d}_i$. The processor computes the route from the position
in the generating report to $\mathbf{g}_i$, using the static obstacle map.
While awaiting another guidance message, robot $i$ consumes waypoints it has
already reached and retains the remaining absolute route. On late delivery,
it also removes waypoints reached since that route's source timestamp.
It neither recomputes A* locally nor translates old relative offsets as if
they referred to its current position. The reported distance remains the
received value until another guidance update arrives.

Let $\widehat{\mathbf{w}}_{i,k}(n)$ denote the $k$th remaining waypoint.
The waypoint features use current self-position,
\begin{equation}
 \widehat{\mathbf{u}}_{i,k}(n)=
 \operatorname{clip}_{[-1,1]}
 \!\left(\frac{\widehat{\mathbf{w}}_{i,k}(n)-\mathbf{q}_i(n)}{h}\right),
 \label{eq:panama-waypoint}
\end{equation}
where clipping is componentwise, $h=7$ is the PANAMA field-of-view radius,
and $k\in\{1,\ldots,K\}$ with $K=5$. Unavailable waypoints contribute zero
features. The retained task vector concatenates normalized reported distance,
the local ever-reached-goal flag $f_i^g(n)$, and these waypoint pairs:
\begin{equation}
 \widehat{\mathbf{v}}_i(n)=
 (\widehat{d}_i/d_i^0,\ f_i^g(n),\
  \widehat{\mathbf{u}}_{i,1},\ldots,\widehat{\mathbf{u}}_{i,K}),
 \label{eq:panama-vector}
\end{equation}
with distance clipped to $[0,1]$ and the initial briefing distance $d_i^0$
as normalizer. A nonpositive or nonfinite normalizer uses the map's larger
dimension, as in the retained task implementation. The goal remains known
even when the route or its reported distance is stale.

Let $\widehat{\mathcal{I}}_{j\rightarrow i}(n)$ be the last received absolute
path segment of peer $j$, and let $\widehat{\rho}_{j\rightarrow i}(n)$ be its
received priority, with smaller values indicating higher priority. Robot $i$
forms the intention channel from peers whose received priority precedes its
own received priority, subject to PANAMA's communication range, local
field of view, and received completion flag. A cached path cell
$\mathbf{w}$ is placed relative to the current $\mathbf{q}_i(n)$.
Writing $\mathcal{R}_i(n)$ for the eligible peers within PANAMA's communication
range and $\widehat{f}_{j\rightarrow i}^{g}$ for their received completion
flags, the contributing set is
\begin{equation}
 \mathcal{H}_i(n)=\left\{j\in\mathcal{R}_i(n)\setminus\{i\}:
 \begin{array}{l}
  \widehat{\rho}_{j\rightarrow i}<\widehat{\rho}_{i\rightarrow i},\\
  \widehat{f}_{j\rightarrow i}^{g}=0
 \end{array}\right\}.
 \label{eq:panama-peer-set}
\end{equation}
For a relative field-of-view cell $\mathbf{z}$, the intention-channel value is
\begin{equation}
 \widehat{O}_i^I(\mathbf{z},n)=\mathbf{1}\!\left\{
 \mathbf{q}_i(n)+\mathbf{z}\in
 \bigcup_{j\in\mathcal{H}_i(n)}\widehat{\mathcal{I}}_{j\rightarrow i}(n)
 \right\}.
 \label{eq:panama-intention-mask}
\end{equation}
Communication loss therefore leaves the peer's old segment in the observation,
while local sensing can still reveal its actual nearby position. Denoting
fresh obstacle and robot channels by $O_i^{\rm obs}(n)$ and
$O_i^{\rm rob}(n)$, the policy input is
\begin{equation}
 \widehat{o}_i(n)=
 \big([O_i^{\rm obs},O_i^{\rm rob},\widehat{O}_i^I],
       \widehat{\mathbf{v}}_i\big).
 \label{eq:panama-observation}
\end{equation}
Each robot continues choosing among north, south, east, west, and wait using
this available input. Communication unavailability does not itself mandate
waiting or insert an external collision-recovery controller.

\subsection{Received Experiences and Continued Learner Operation}
After joint action $n$, robot $i$ produces the complete transition
\begin{equation}
 \zeta_i(n)=
 (\widehat{o}_i(n),a_i(n),r_i(n+1),
  \widehat{o}_i(n+1),z(n+1)),
 \label{eq:panama-transition}
\end{equation}
where $z$ is the task simulator's terminal flag. The tuple is captured at
$t_{n+1}$ together with robot, episode, joint-step, and deployed-version
identifiers. Upload delay changes only its arrival time,
\begin{equation}
 b_i^E(n)=t_{n+1}+\delta_i^E(n).
 \label{eq:panama-upload}
\end{equation}
The transition's observations, action, reward, and terminal flag remain tied
to their generating step, regardless of later robot motion or weight updates.

At learner time $t$, the received subset of the $N$ expected transitions
from joint action $n$ is
\begin{equation}
 \mathcal{U}_n(t)=\{i\in\mathcal{N}:b_i^E(n)\leq t\}.
 \label{eq:panama-subset}
\end{equation}
The replay buffer contains only fully received transitions, subject to its
capacity and curriculum-phase retention rule. In particular,
\begin{equation}
 \mathcal{B}(t)\subseteq
 \bigcup_{n:\,t_{n+1}\leq t}
 \{\zeta_i(n):i\in\mathcal{U}_n(t)\}.
 \label{eq:panama-replay}
\end{equation}
The set on the right ranges over all generated joint actions. The learner
does not require $\mathcal{U}_n(t)=\mathcal{N}$ before learning. It can use
received experiences from unaffected robots and previously delivered replay
samples while other uploads remain pending. If all current uploads are
delayed, existing replay still supports updates. A learning opportunity is
skipped only when insufficient replay data are available. Late experiences
from a completed curriculum phase are excluded from the new phase's replay.

For each sampled transition $\zeta=(o,a,r,o',z)$, the retained Double DQN
target and weighted Huber loss are
\begin{align}
 a_\zeta^*&=\operatorname*{arg\,max}_{a'}Q_\theta(o',a'),\\
 y_\zeta&=r+\gamma(1-z)Q_{\theta^-}(o',a_\zeta^*),\\
 \mathcal{L}(\theta)&=\frac{1}{M}
       \sum_{\zeta\in\mathcal{J}(t)}
       w_\zeta\,\ell_{\rm H}\big(Q_\theta(o,a)-y_\zeta\big),
 \label{eq:panama-loss}
\end{align}
where $\mathcal{J}(t)$ is a PER minibatch from $\mathcal{B}(t)$ and
$w_\zeta$ is its importance-sampling weight. AdamW updates $\theta$ and the
target follows
$\theta^-\leftarrow(1-\eta_{\rm tar})\theta^-+\eta_{\rm tar}\theta$.
The delayed upload can change data availability and replay age, but does not
replace the learning rule with an update based on incomplete tuples.
The learner is off-policy: an experience collected with an older deployed
version remains a complete training sample, and the learner evaluates its
target with the current policy and target networks. Recording the generating
version permits analysis of this mismatch without requiring every robot to
hold the latest model before a learner update.

\subsection{Complete Model Deployment and Per-Robot Staleness}
Let $\Theta^\ell$ be the immutable complete model snapshot published at time
$u_\ell$, with increasing publication version $\ell$. For robot $i$, its
completed downlink arrival is
\begin{equation}
 b_i^W(\ell)=u_\ell+\delta_i^W(\ell).
 \label{eq:panama-download}
\end{equation}
Immediately before action $n$, robot $i$ uses
\begin{align}
 v_i(n)&=\max\{\ell:b_i^W(\ell)\leq t_n\},\\
 \widetilde{\theta}_i(n)&=\Theta^{v_i(n)},
 \label{eq:panama-deployed}
\end{align}
with a complete initial model available at $t_0$. A receiver installs a newer
snapshot atomically after full reception, before its next action. It never
mixes layers, partially received parameters, or components of two versions.
A late snapshot with $\ell\leq v_i(n)$ cannot roll the deployed model back.

The action is selected by the standard epsilon-greedy rule,
\begin{align}
 a_i^*(n)&=\operatorname*{arg\,max}_{a\in\mathcal{A}}
     Q_{\widetilde{\theta}_i(n)}(\widehat{o}_i(n),a),\\
 \pi_i(a\mid\widehat{o}_i;n)&=\frac{\epsilon_n}{|\mathcal{A}|}
       +(1-\epsilon_n)\mathbf{1}\{a=a_i^*(n)\}.
 \label{eq:panama-action}
\end{align}
Exploration retains PANAMA's schedule and depends on generated environment
steps, not on delivered experience counts. A missing downlink update leaves
$\widetilde{\theta}_i(n)$ unchanged while the robot continues acting.
Other robots can install newer snapshots independently, and the centralized
learner can continue optimizing its current model. Publications occur every
$H$ generated joint actions, with the same $H$ in all learning conditions.

Let $\ell_{\rm pub}(t_n)$ be the latest published version. We separately
measure model-version lag and model age,
\begin{align}
 L_i^W(n)&=\ell_{\rm pub}(t_n)-v_i(n),\\
 A_i^W(n)&=t_n-u_{v_i(n)}.
 \label{eq:panama-model-age}
\end{align}
Experience delivery age is $b_i^E(n)-t_{n+1}$; replay age at sampling time
$t$ is $t-t_{n+1}$. These quantities distinguish old policy execution from
old learning data without treating either as a halted mission.

\subsection{Operation and Separation of Communication Effects}
Algorithm~\ref{alg:panama-delay} summarizes the event-driven execution.
The flag $\chi_L$ equals one when learning is enabled and zero during
frozen-policy inference.
One learning opportunity follows each generated joint action. Generated
steps, delivered experiences, completed learner updates, and published models
are counted independently. Delayed uploads cannot redefine the step counter
or prevent other robots from contributing experiences. The simulator advances
its communication clock with robot actions and processes arrivals at their
timestamps, without a wall-clock sleep representing communication delay.

\begin{algorithm}[t]
\color{panamaaddition}
\captionsetup{font={color=panamaaddition},labelfont={color=panamaaddition}}
\caption{PANAMA with timestamped information and learning transport}
\label{alg:panama-delay}
\begin{algorithmic}[1]
\Require Initial caches, model $\Theta^0$, replay $\mathcal{B}$, $\chi_L$
\For{each generated joint action $n$}
  \State Process completed arrivals up to decision time $t_n$
  \State Accept newer information and complete model versions
  \For{each robot $i\in\mathcal{N}$}
    \State Sense locally and consume reached cached waypoints
    \State Construct $\widehat{o}_i(n)$ and act with $\Theta^{v_i(n)}$
  \EndFor
  \State Apply joint action using the true physical transition
  \State Advance event time and process completed arrivals
  \For{each robot $i\in\mathcal{N}$}
    \State Submit timestamped mission reports to the transport
    \State Capture and queue the complete $\zeta_i(n)$
  \EndFor
  \State Insert newly received eligible experiences into $\mathcal{B}$
  \If{$\chi_L=1$ and $|\mathcal{B}|\geq M$}
    \State Perform one standard PER/Double DQN update
  \EndIf
  \If{$\chi_L=1$ and $n+1$ is a multiple of $H$}
    \State Publish and queue complete snapshot $\Theta^{\ell+1}$
  \EndIf
\EndFor
\end{algorithmic}
\end{algorithm}

For frozen-policy mission analysis, all robots use one fixed snapshot and
learning is disabled. The transport changes only the selected guidance,
intention, or priority streams; unselected inputs remain current. For
learning-transport analysis, mission information remains current and we
compare immediate upload and deployment with delayed upload, delayed
deployment, and their combination. The same initialization and nominal PANAMA
parameters apply to those learning conditions. This separation attributes
mission effects to stale communicated information and learning effects to
received-data or deployed-model staleness. It makes no assumption that a
communication delay necessarily reduces success.
\end{panamaaddition}
% END PANAMA ADDITIONS: PANAMA under communication delay
